\chapter{BRAVA-GNN}
\section{Overview}

BRAVA-GNN (Betweenness Ranking Approximation via Degree Mass GNN) is a lightweight graph neural network architecture designed to approximate betweenness centrality rankings for both directed and undirected networks. Its primary contribution is achieving high ranking accuracy while maintaining a parameter count that is strictly independent of the input graph size.

\section{Key Design Principles}

\subsection{Size-Invariant Node Features}
Instead of relying on node-specific embeddings that scale with the number of nodes, BRAVA-GNN uses \textbf{multi-hop degree mass} as its input feature. The degree mass captures the cumulative connectivity around a node up to a certain distance. Using degree masses up to order five produces a fixed-dimensional input vector, significantly reducing the model size.

\subsection{Dual Message Passing}
To effectively model directed networks where shortest paths are asymmetric, BRAVA-GNN introduces a dual message-passing mechanism. It applies two parallel streams:
\begin{itemize}
    \item An \textbf{incoming stream} over the adjacency matrix $A$.
    \item An \textbf{outgoing stream} over the transposed adjacency matrix $A^T$.
\end{itemize}
These two streams share weights to keep the architecture compact. The intermediate representations from both directions are then fused multiplicatively, effectively rewarding nodes that are reachable from many sources and can reach many destinations.

\subsection{Mixed Synthetic Training Distribution}
To ensure the model generalizes to diverse real-world topologies without needing domain-specific retraining, BRAVA-GNN is trained entirely on a dataset of synthetic graphs. In the original paper, comprehensive experiments and ablation studies revealed that the best performance is achieved by using only Scale-Free (SF) graphs for training. Because these specific graphs yielded the optimal results in the original evaluation, we exclusively use directed and undirected Scale-Free graphs for our training process.

\subsection{Training Configuration}
We strictly follow the optimal training configuration identified in the original paper. The architecture uses $L=2$ message-passing layers with a hidden dimension of $n_{hid} = 12$. The final prediction MLP consists of three linear layers (dimensions $12 \to 24 \to 24 \to 1$) with ReLU activations and a dropout rate of $p=0.3$. Row-wise $L_2$ normalization is applied after each message-passing step to improve stability. The model is trained for 10 epochs using the Adam optimizer with a learning rate of $5 \times 10^{-3}$.

\section{Performance}

By relying on multi-hop structural features and a compact dual-stream architecture, BRAVA-GNN uses approximately 1,300 parameters—orders of magnitude fewer than previous GNN-based approaches. Despite its small size, it achieves state-of-the-art Kendall $\tau$ correlation on real-world test networks, and offers competitive inference latency, achieving massive speedups over exact algorithms.
